\begin{abstract}
Multi-agent language-model systems usually optimize the final aggregate while leaving the behavior that produced it unobserved. This is problematic when a complete, independently reasoned report is more costly than a short answer or imitation. We introduce a \emph{game-driven observable semantic-consensus mechanism} that uses Assumption--Evidence--Inference--Conclusion (AEIC) reports as a parseable interface between a mechanism designer and LLM agents. Before solving, each agent observes a public credit--cost rule and chooses among low-effort reporting, structured reasoning, verification, and peer imitation. The chosen action controls the actual generation procedure and is logged before settlement, enabling a direct follower-response test rather than retrospective response scoring. AEIC is treated as an observation and settlement interface, not assumed to be an intrinsically accuracy-improving prompt template.

Across 450 paired test questions from StrategyQA, MMLU-Hard, and TruthfulQA, a static process contract over the AEIC interface increases high-effort choices from 0.46\% to 56.93\%; under the joint contract, execution-prompt, and tiered-resource configuration, majority-vote accuracy rises from 68.00\% to 82.22\% ($+14.22$ percentage points; Holm-adjusted McNemar $p=1.23\times10^{-11}$). An outcome-only reward raises effort only to 3.08\% and does not improve accuracy. A matched-compute control does not reproduce the contract's behavioral response; in a separate interface test, AEIC outperforms free text but not generic structure. Reward and cost sweeps produce the predicted effort response, field ablations do not establish the necessity of any single AEIC field, and a held-out GLM backend reproduces neither gain. Eight frozen suites plus two supplementary experiments contain 6,360 formal question-level records. The supported claim is bounded: public process contracts can change follower resource actions, while the observed quality gain belongs to the joint contract-and-execution configuration rather than to AEIC formatting alone.

\keywords{Multi-agent LLM systems \and Mechanism design \and Structured reasoning \and Deliberation \and Follower response}
\end{abstract}
